\section{Ipergeometrica}
La distribuzione ipergeometrica descrive il numero di successi ottenuti estraendo senza reinserimento un campione di dimensione $n$ da una popolazione finita di $N$ elementi, di cui $K$ sono classificati come successi e $N-K$ come insuccessi.
\[
\mathcal{H}(N,K,n)
\]

\begin{figure}[H]
\centering

\begin{minipage}{0.45\textwidth}
\centering
\begin{tikzpicture}
\begin{axis}[
    axis lines = middle,
    ylabel = {$\mathbb{P}(X=k)$},
    xlabel = {$k$},
    width = 0.90\textwidth,
    height = 0.90\textwidth,
    ymin = 0, ymax = 0.18,
    xmin = -1, xmax = 31,
    legend style={at={(axis cs: 1,0.17)}, anchor=north west, font=\tiny},
    legend cell align=left,
]

% --- Hypergeom(N=100,K=40,n=30) ---
\addplot[
    only marks,
    black,
    mark=*,
    mark options={fill=black, scale=0.7}
]
coordinates {
(5,0.001)(6,0.004)(7,0.011)(8,0.024)(9,0.044)
(10,0.069)(11,0.098)(12,0.125)(13,0.142)(14,0.145)
(15,0.134)(16,0.111)(17,0.082)(18,0.053)
(19,0.030)(20,0.015)(21,0.006)(22,0.002)
};

% --- Hypergeom(N=100,K=20,n=30) ---
\addplot[
    only marks,
    gray!60,
    mark=*,
    mark options={fill=gray!60, scale=0.7}
]
coordinates {
(0,0.003)(1,0.017)(2,0.048)(3,0.092)(4,0.132)
(5,0.155)(6,0.157)(7,0.139)(8,0.108)(9,0.074)
(10,0.045)(11,0.023)(12,0.010)(13,0.004)(14,0.001)
};

\end{axis}
\end{tikzpicture}
\end{minipage}
\hfill
\begin{minipage}{0.45\textwidth}
\centering

\begin{tikzpicture}
\begin{axis}[
    axis lines = middle,
    ylabel = {$F(k)$},
    xlabel = {$k$},
    width = 0.90\textwidth,
    height = 0.90\textwidth,
    ymin = 0, ymax = 1.1,
    xmin = -1, xmax = 31,
    legend style={at={(axis cs: 15,0.25)}, anchor=north west, font=\tiny},
    legend cell align=left,
]

% --- CDF Hypergeom(N=100,K=40,n=30) ---
\addplot[
    jump mark left,
    black,
    mark=*,
    mark options={fill=black, scale=0.7}
]
coordinates{
(0,4.02639294e-09)
(1,1.59886765e-07)
(2,2.91423177e-06)
(3,3.25164852e-05)
(4,0.00024996245)
(5,0.00141298772)
(6,0.00612431695)
(7,0.02096773258)
(8,0.05802744464)
(9,0.13235803521)
(10,0.25333107136)
(11,0.41427080904)
(12,0.59021881590)
(13,0.74885530151)
(14,0.86706008543)
(15,0.93990925597)
(16,0.97702670020)
(17,0.99263553782)
(18,0.99803720732)
(19,0.99956893745)
(20,0.99992276711)
(21,0.99998884174)
(22,0.99999871828)
(23,0.99999988499)
(24,0.99999999212)
(25,0.99999999960)
(26,0.99999999999)
(27,1.00000000000)
(28,1.00000000000)
(29,1.00000000000)
(30,1.00000000000)
};

\addlegendentry{$N=100,\,K=40,\,n=30$}

% --- CDF Hypergeom(N=100,K=20,n=30) ---
\addplot[
    jump mark left,
    gray!60,
    mark=*,
    mark options={fill=gray!60, scale=0.7}
]
coordinates{
(0,0.003)(1,0.020)(2,0.068)(3,0.160)(4,0.292)
(5,0.447)(6,0.604)(7,0.743)(8,0.851)
(9,0.925)(10,0.970)(11,0.993)(12,0.999)
(13,1.000)
};
\addlegendentry{$N=100,\,K=20,\,n=30$}

\end{axis}
\end{tikzpicture}

\end{minipage}
\end{figure}

\bigskip
\begin{table}[H]
\centering
\renewcommand{\arraystretch}{2}
\begin{tabular}{@{}ll@{}}
\toprule
\textbf{Supporto} 
& $\{\max(0,n-(N-K)),\dots,\min(n,K)\}\subset\mathbb{Z}$ \\

\midrule

\textbf{Funzione di Densità} 
& $\displaystyle
\mathbb{P}(X=k)=
\frac{\binom{K}{k}\binom{N-K}{\,n-k\,}}{\binom{N}{n}}
$ \\

\midrule

\textbf{Funzione di Ripartizione}
& $F(k)=\displaystyle\sum_{j=0}^{\lfloor k\rfloor}
\frac{\binom{K}{j}\binom{N-K}{\,n-j\,}}{\binom{N}{n}}$ \\

\midrule

\textbf{Valore atteso} 
& $\mathbb{E}[X]=n\frac{K}{N}$ \\

\midrule

\textbf{Varianza}
& $Var(X)=n\frac{K}{N}\left(1-\frac{K}{N}\right)
\frac{N-n}{N-1}$ \\

\bottomrule
\end{tabular}
\end{table}
